% Blueprint for the formalization of Freedman's disc embedding theorem.
% Source: S. Behrens, B. Kalmár, M. H. Kim, M. Powell, A. Ray (eds.),
%   The Disc Embedding Theorem, Oxford University Press, 2021. Cited as [DET].
% Secondary: M. H. Freedman, "Bing Topology and Casson Handles",
%   2013 lectures, notes by S. Behrens (freely available).
%
% Conventions: \lean{...} names the Lean declaration, \leanok marks a
% statement (or proof) that is formalized, \uses{...} records dependencies.
% Chapter numbers refer to [DET].

\chapter{Overview and conventions}

The goal is a complete Lean formalization of the disc embedding theorem
[DET, Ch.~11] and its headline consequence, the topological
4-dimensional Poincar\'e conjecture. We follow the book chapter by chapter.

\paragraph{Models.} We model the closed $n$-disc $D^n$ as the closed unit ball
of $\mathbb{R}^n$ with the sup metric (the cube $[-1,1]^n$) and $S^{n-1}$ as its
boundary. These are homeomorphic to the round disc and sphere; every statement
in Part I is invariant under homeomorphism.

\begin{definition}[Model disc and sphere]
  \label{def:disc}
  \lean{DiscEmbedding.Disc, DiscEmbedding.Sphere}
  \leanok
  $D^n = \{x \in \mathbb{R}^n : \|x\|_\infty \le 1\}$ and
  $S^{n-1} = \partial D^n$. In Lean, \texttt{Sphere (n+1)} is the $n$-sphere.
\end{definition}

\begin{definition}[Topological ball]
  \label{def:top-ball}
  \lean{DiscEmbedding.IsTopBall}
  \leanok
  \uses{def:disc}
  A subset $B \subseteq X$ is a topological $n$-ball if it is homeomorphic,
  with the subspace topology, to $D^n$.
\end{definition}

% ------------------------------------------------------------------
\chapter{Part I: Decomposition space theory}

\section{The Schoenflies theorem after Mazur, Morse and Brown [DET, Ch.~3]}

\begin{definition}[Cellular set]
  \label{def:cellular}
  \lean{DiscEmbedding.IsCellular}
  \leanok
  \uses{def:disc}
  $A \subseteq M$ is cellular if $A = \bigcap_i e_i(D^n)$, where each
  $e_i\colon \mathbb{R}^n \to M$ is an open embedding and
  $e_{i+1}(D^n) \subseteq e_i(\operatorname{int} D^n)$. (In a manifold this is
  equivalent, via invariance of domain, to asking for arbitrary $n$-balls $B_i$
  with $B_{i+1}\subseteq \operatorname{int} B_i$.)
\end{definition}

\begin{lemma}[Radial squeeze]
  \label{lem:squeeze}
  \lean{DiscEmbedding.squeezeHomeomorph}
  \leanok
  For $0<\delta<r<1$ there is a homeomorphism of $\mathbb{R}^n$ that is the
  identity outside $D^n$, preserves $D^n$, and maps the closed $r$-ball into the
  closed $\delta$-ball.
\end{lemma}

\begin{proof}
  \leanok
  Rescale each vector radially by the piecewise-linear order automorphism of
  $\mathbb{R}$ through $(0,0)$, $(r,\delta)$, $(1,1)$ with slope $1$ beyond $1$.
\end{proof}

\begin{lemma}[Transplanting]
  \label{lem:transplant}
  \lean{DiscEmbedding.transplant}
  \leanok
  If $e\colon E\to M$ is an open embedding into a Hausdorff space and $\sigma$ is a
  homeomorphism of $E$ that is the identity outside a compact set, then
  $e\sigma e^{-1}$ on $e(E)$, extended by the identity, is a homeomorphism of $M$.
\end{lemma}

\begin{proof}
  \leanok
  Continuity follows on the open set $e(E)$ and on the open complement of the
  compact set $e(K)$, which together cover $M$.
\end{proof}

\begin{theorem}[Cellular sets shrink]
  \label{thm:cellular-shrink}
  \lean{DiscEmbedding.IsCellular.satisfiesBSC_of_collapse, DiscEmbedding.cellular_quotient_approximable}
  \leanok
  \uses{def:cellular, def:approx, def:bsc}
  If $A \subseteq M$ is cellular and $q\colon M \to Y$ collapses exactly $A$ to a
  point, then $q$ satisfies the shrinking criterion. Hence if $M$ is compact, $q$ is
  approximable by homeomorphisms; in particular $S^n/A \cong S^n$.
\end{theorem}

\begin{proof}
  \uses{thm:bsc, lem:squeeze, lem:transplant}
  \leanok
  Given $\varepsilon$, choose $i$ with $q(e_i(D^n))$ within $\varepsilon/2$ of the
  collapse point. Then $A\subseteq e_{i+1}(D^n)$ lies in $e_i$ of a closed
  $r$-ball with $r<1$. Transplant the radial squeeze through $e_i$, with $\delta$
  chosen so that $e_i$ maps the $\delta$-ball into an $\varepsilon/3$-ball. The
  result moves points only within $e_i(D^n)$ and shrinks $A$ to diameter
  $<\varepsilon$; every other fiber is a point.
\end{proof}

\subsection*{IOUs from algebraic topology}

The following classical results are used but not yet formalized. They are stated
in \texttt{IOU.lean} with \texttt{sorry}.

\begin{theorem}[IOU: no-retraction]
  \label{iou:noretract}
  \lean{DiscEmbedding.no_retraction}
  \leanok
  There is no retraction of $D^n$ onto $\partial D^n$.
\end{theorem}

\begin{theorem}[IOU: invariance of domain]
  \label{iou:iod}
  \lean{DiscEmbedding.invariance_of_domain}
  \leanok
  An injective continuous map from an open subset of $\mathbb{R}^n$ to
  $\mathbb{R}^n$ has open image.
\end{theorem}

\begin{theorem}[IOU: Jordan--Brouwer separation]
  \label{iou:jb}
  \lean{DiscEmbedding.jordan_brouwer_separation}
  \leanok
  The complement of an embedded $S^{n-1}$ in $S^n$ has exactly two components,
  each with the sphere as frontier.
\end{theorem}

\subsection*{Brown's proof}

We follow M.~Brown, \emph{A proof of the generalized Schoenflies theorem}, Bull.\ AMS
66 (1960). An \emph{inverse set} of a map is a fiber with at least two points.

\begin{theorem}[Brown, Thm.~0]
  \label{thm:brown0}
  \uses{iou:noretract, iou:iod, iou:jb}
  If $Q$ is an $n$-cell and $f\colon Q\to S^n$ has finitely many inverse sets, all
  in $Q$, then $f(Q)$ is the union of $f(\partial Q)$ and one complementary domain of
  $f(\partial Q)$.
\end{theorem}

\begin{theorem}[Brown, Thm.~1]
  \label{thm:brown1}
  \lean{DiscEmbedding.brown_collapse_cellular}
  \leanok
  \uses{def:cellular, lem:squeeze, lem:transplant}
  A cellular subset $M$ of the open cube admits a continuous surjection of
  $\mathbb{R}^n$, the identity outside the open cube, whose only inverse set is $M$.
\end{theorem}

\begin{proof}
  \leanok
  Set $h_0=\mathrm{id}$ and $h_{k+1}=h_k\tau_k$, where $\tau_k$ is supported in $B_k$,
  maps $B_k$ into itself, and makes $\operatorname{diam} h_{k+1}(B_{k+1})<1/(k+1)$.
  Then $h_j(B_k)\subseteq h_k(B_k)$ for $j\ge k$, so the $h_k$ converge uniformly. The
  limit collapses $M$ and is injective elsewhere, and it is surjective as a uniform
  limit of surjections.
\end{proof}

\begin{theorem}[Brown, Thm.~2]
  \label{thm:brown2}
  \uses{thm:brown1}
  Let $D$ be a complementary domain of an $(n-1)$-sphere in $S^n$. If $f$ maps
  $\overline{D}$ onto an $n$-cell with a single, cellular inverse set $M\subseteq D$,
  then $\overline{D}$ is an $n$-cell.
\end{theorem}

\begin{theorem}[Brown, Thm.~3]
  \label{thm:brown3}
  \uses{thm:brown0, lem:transplant}
  If $f\colon Q\to S^n$ has a single inverse set $M\subseteq\operatorname{int}Q$,
  then $M$ is cellular.
\end{theorem}

\begin{theorem}[Brown, Thm.~4]
  \label{thm:brown4}
  \lean{DiscEmbedding.brown_two_inverse_sets}
  \leanok
  \uses{thm:brown0, thm:brown3}
  If $f\colon S^n\to S^n$ is surjective with exactly two inverse sets, both are
  cellular.
\end{theorem}

\begin{theorem}[Brown, Thm.~5]
  \label{thm:brown5}
  \lean{DiscEmbedding.brown_closure_isTopBall}
  \leanok
  \uses{thm:brown4, thm:brown2}
  The closure of each complementary component of a bicollared $(n-1)$-sphere in
  $S^n$ is an $n$-ball.
\end{theorem}

\begin{theorem}[Generalized Schoenflies]
  \label{thm:schoenflies}
  \lean{DiscEmbedding.generalized_schoenflies}
  \leanok
  \uses{def:top-ball}
  If $e\colon S^{n-1}\times[-1,1] \hookrightarrow S^n$ is an embedding, then
  $S^n \setminus e(S^{n-1}\times 0)$ has two components and the closure of each
  is an $n$-ball.
\end{theorem}

\begin{proof}
  \uses{iou:jb, thm:brown5}
  \leanok
  Jordan--Brouwer gives the two components; Brown's Theorem~5 shows each closure
  is a ball.
\end{proof}

\section{Decomposition space theory and the Bing shrinking criterion [DET, Ch.~4]}

\begin{definition}[Approximable by homeomorphisms]
  \label{def:approx}
  \lean{DiscEmbedding.ApproximableByHomeomorphisms}
  \leanok
  $f\colon X \to Y$ is approximable by homeomorphisms if for every
  $\varepsilon>0$ there is a homeomorphism $g\colon X\to Y$ with
  $d(g(x), f(x))<\varepsilon$ for all $x$.
\end{definition}

\begin{definition}[Shrinking criterion]
  \label{def:bsc}
  \lean{DiscEmbedding.SatisfiesBSC}
  \leanok
  $f\colon X\to Y$ satisfies the Bing shrinking criterion if for every
  $\varepsilon>0$ there is a homeomorphism $h\colon X\to X$ with
  $d(fh(x), f(x)) < \varepsilon$ for all $x$ and
  $\operatorname{diam} h(f^{-1}(y)) < \varepsilon$ for all $y$.
\end{definition}

\begin{lemma}[Easy direction]
  \label{lem:bsc-easy}
  \lean{DiscEmbedding.ApproximableByHomeomorphisms.satisfiesBSC}
  \leanok
  \uses{def:approx, def:bsc}
  If $Y$ is compact and $f$ is approximable by homeomorphisms, then $f$
  satisfies the shrinking criterion.
\end{lemma}

\begin{proof}
  \leanok
  Take $g$ within $\varepsilon/2$ of $f$; by uniform continuity of $g^{-1}$ pick
  $\eta$; take $g'$ within $\min(\eta/2,\varepsilon/2)$ of $f$; set
  $h=g^{-1}g'$.
\end{proof}

\begin{theorem}[Bing shrinking criterion]
  \label{thm:bsc}
  \lean{DiscEmbedding.bing_shrinking_criterion}
  \leanok
  \uses{def:approx, def:bsc}
  A continuous surjection $f\colon X\to Y$ between compact metric spaces is
  approximable by homeomorphisms if and only if it satisfies the shrinking
  criterion.
\end{theorem}

\begin{proof}
  \uses{lem:bsc-easy}
  \leanok
  The forward direction is Lemma~\ref{lem:bsc-easy}. For the converse, let
  $A$ be the closure of $\{f\circ h : h\in\mathrm{Homeo}(X)\}$ in $C(X,Y)$ with
  the sup metric; it is complete, hence Baire. Every map in $A$ is surjective.
  For each $n$, the maps that separate points at distance $\ge 1/(n+1)$ form an
  open subset of $A$ (by compactness of $X\times X$), and the shrinking criterion,
  together with uniform continuity of $h^{-1}$, shows it is dense. A map in the
  intersection that is within $\varepsilon$ of $f$ is a continuous bijection
  from a compact space to a Hausdorff space, hence a homeomorphism.
\end{proof}

\section{Remaining chapters of Part I}

These are not yet stated in Lean. Each needs an explicit construction of a
decomposition, which is where the formalization work starts to become
geometric.

\begin{itemize}
  \item \textbf{Ch.~5} The Alexander gored ball and the Bing decomposition:
        the Bing decomposition of $S^3$ shrinks.
  \item \textbf{Ch.~6} A decomposition that does not shrink (Bing-2).
  \item \textbf{Ch.~7} The Whitehead decomposition.
  \item \textbf{Ch.~8} Mixed Bing--Whitehead decompositions.
  \item \textbf{Ch.~9} Shrinking starlike sets.
  \item \textbf{Ch.~10} The ball to ball theorem. \emph{This is the Part I
        result used later in the proof; its hypotheses must be transcribed
        exactly from [DET] before it is stated in Lean.}
\end{itemize}

% ------------------------------------------------------------------
\chapter{Parts II--IV: Building skyscrapers, and skyscrapers are standard}

Not yet started. Infrastructure needed (none of it is in Mathlib):
topological 4-manifolds with boundary and handle decompositions; generic
immersions of discs and spheres; intersection and self-intersection numbers
(with $\mathbb{Z}[\pi_1]$ coefficients); Whitney and finger moves; gropes,
towers and skyscrapers [DET, Ch.~12--14]; the ``picture camp'' calculus
[DET, Ch.~13]; the collar adding lemma [DET, Ch.~25].

\begin{theorem}[Topological 4-dimensional Poincar\'e conjecture]
  \label{thm:poincare4}
  \lean{DiscEmbedding.poincare_four}
  \leanok
  \uses{def:disc}
  A closed topological 4-manifold homotopy equivalent to $S^4$ is homeomorphic
  to $S^4$.
\end{theorem}
